% Generated draft: requires mathematical and source review.
\documentclass[11pt]{article}
\usepackage[margin=25mm]{geometry}
\usepackage{fontspec}
\setmainfont{Latin Modern Roman}
\usepackage{amsmath,amssymb,mathtools,unicode-math}
\setmathfont{Latin Modern Math}
\usepackage{enumitem,needspace,xurl,hyperref}
\setlength{\parindent}{0pt}
\setlength{\parskip}{5pt}
\setlength{\emergencystretch}{3em}
\newcommand{\N}{\mathbb{N}}
\newcommand{\Z}{\mathbb{Z}}
\newcommand{\Q}{\mathbb{Q}}
\newcommand{\R}{\mathbb{R}}
\newcommand{\C}{\mathbb{C}}
\newcommand{\F}{\mathbb{F}}
\newcommand{\A}{\mathbb{A}}
\newcommand{\PP}{\mathbb{P}}
\newcommand{\E}{\mathbb{E}}
\newcommand{\eps}{\varepsilon}
\newcommand{\dd}{\,\mathrm{d}}
\begin{document}
\section*{UnsolvedMath ID: 6200069}
\textbf{Source problem code:} catalogue entry
\subsection{Definitions and mathematical statement}
For a metric space $(X,d)$ with basepoint $x_0$, normalized distance functions are
\[
 h_y(x)=d(x,y)-d(x_0,y).
\]
A standard horofunction is an appropriate limit of these functions as $y$ escapes, normalized by $h(x_0)=0$; for proper spaces one commonly uses uniform convergence on compact sets. Whether the source allows all metric-function limits or only boundary points is not stated.

Let $x_n=x_n(\omega)$ be the intended random path and, when the limit is finite and exists, write
\[
 A(\omega)=\lim_{n\to\infty}\frac1n d(x_0,x_n(\omega)).
\]
The source gives neither the probability law nor the metric/cocycle assumptions ensuring this drift. It does not say whether ergodicity makes $A$ constant.

\textbf{Conjecture attributed to Anders Karlsson.} Almost surely there exists one horofunction $h=h_\omega$ such that
\[
 \lim_{n\to\infty}-\frac1n h(x_n(\omega))=A(\omega).
\]
The horofunction may depend on the path; no measurable-selection requirement is explicitly supplied.

The record separately reports Karlsson's approximation theorem with quantifiers
\[
 \forall\varepsilon>0\ \exists h_\varepsilon\ \exists N(\varepsilon)\ 
 \forall n\geq N(\varepsilon):\qquad
 A-\varepsilon\leq-\frac1n h_\varepsilon(x_n)\leq A+\varepsilon.
\]
Here $h_\varepsilon$ is fixed for that tolerance and the whole subsequent tail, but may change with $\varepsilon$. This is not silently identified with one $h$ satisfying the exact limit. The hypotheses and probability scope of the cited approximation theorem are absent.

The research summary marks the exact statement solved and invokes horofunction drift, Kingman-style arguments, and Karlsson--Margulis. It names Karlsson's ``Non-reversibility and equivalent conditions for convexity'' without a precise matching theorem. These attributed claims do not authenticate the missing path model, drift conditions, or exact result in this excerpt. Neighboring local registry entries were checked and do not supply that context.
\subsection{Short English statement}
For the intended random path with linear escape rate $A$, find one horofunction whose negative values along the path grow at rate $A$ almost surely. The source already reports an epsilon-dependent approximation, which is weaker in its quantifiers. The path law and hypotheses behind the claimed exact theorem are missing.
\subsection{Sources}
\begin{itemize}
\item Local UnsolvedMath registry, permanent ID 6200069, source code catalogue entry, ``Boundaries of Groups and Kleinian Groups --- Problem 69''; attribution field: Conjecture of Anders Karlsson; proposed year 2005. Status solved, classification SOLVED-IN-LITERATURE, updated 2026-08-06.
\item Kapovich, \emph{Problems on Boundaries of Groups and Kleinian Groups} (2005), Problem 69, PDF page 18, \url{https://www.math.ucdavis.edu/~kapovich/EPR/problems.pdf}. Reported PDF-text-normalized extraction and access: 2026-07-29. Current status NEEDS\_REVIEW; the source labels it a problem. Rights: public author-hosted PDF; statement-reuse terms NEEDS\_REVIEW.
\item The statement attributes the tolerance-dependent result to Karlsson. The summary cites Karlsson, ``Non-reversibility and equivalent conditions for convexity,'' and Karlsson--Margulis without complete publication details or theorem hypotheses; Kingman-style arguments are mentioned as context.
\end{itemize}
\end{document}
